\documentclass[aps,pra,reprint,groupedaddress,floatfix]{revtex4-2}
\usepackage{graphicx}
\usepackage{amsmath,amssymb}
\usepackage{placeins}
% Float tuning (reprint two-column carries two large figure*s: D04 2x2, D06 trio).
% Legibility-first: keep panel sizes, relax placement fractions instead.
\renewcommand{\topfraction}{0.9}
\renewcommand{\bottomfraction}{0.8}
\renewcommand{\textfraction}{0.07}
\renewcommand{\floatpagefraction}{0.6}
\usepackage{xeCJK}
\setCJKmainfont{Noto Serif CJK SC}

\graphicspath{{figures/}}

\begin{document}

\title{SSH XXZ 链的模型与归一化拓扑反映标记的有限尺寸性质}

\author{邱雪}\thanks{共同第一作者。}
\author{韩懿杰}\thanks{共同第一作者。}
\author{余金龙}
\affiliation{海南大学物理与光电工程学院，海口 570228}

\date{\today}

\begin{abstract}
变分本征求解器已成为近期量子模拟的基石，但拓扑相在真机上难以认证，
因为理论家偏好的标记很少可直接测量。
本文以保对称 orbit 拟设在相匹配初态上制备二聚化 XXZ 链中竞争的平庸、拓扑、
反铁磁三相。以归一化反射标记 $\tilde{Z}_{\mathcal R}$ 画出 $L=8$ 有限尺寸相图，
分辨三相，能隙标度确认有限尺寸效应可控。
态制备只用可直接测量的标记——AFM 结构因子 $S(\pi)$ 与 string 算符，
优质比特筛选与张量积读出缓解把协议搬上超导真机。
真机与模拟机趋势以各自内部 $0$--$1$ 归一化在同等口径下比较，
覆盖两个二聚化取值与两种可直接测量标记。
\end{abstract}

\maketitle

\section{引言}
\label{sec:intro}

变分量子本征求解器已成为近期量子模拟的基石，用经典优化循环换取电路深度。
但对拓扑相，标准变分流程多一道障碍：理论家偏好的序参量往往在真机上最难测量。

真机上的拓扑相研究沿直接构造与变分优化两条线展开。超导处理器上，
Satzinger 等人~\cite{satzinger2021realizing}制备了 toric code 基态并测得
拓扑纠缠熵，以拓扑序本身为目标；这类直接构造保真度尚可，但只适用于严格可解
定点态。变分方向上，Sun 等人~\cite{sun2023efficient}为对称保护拓扑相设计了
初态层加变分层拟设，并在 IBM 设备上验证了 string 序与边缘模，但拟设结构
很大程度上是对每个目标相试凑出来的。

相匹配初态的主要思想虽简单，其在一般非可积模型上的全部潜力——用单一拟设族
覆盖竞争序、标记经得起向真机迁移——尚未挖掘。
我们借鉴的邻近文献都已成熟：一维对称保护相由射影表示与矩阵乘积态
分类~\cite{chen2011classification,schuch2011classifying}；
随机测量以很少的设置重构纠缠与多体不变量~\cite{huang2020predicting,elben2022toolbox}；
变分算法及其可训练性极限已有综述~\cite{cerezo2021variational,tilly2022variational,mcclean2018barren}；
误差缓解覆盖概率抵消与零噪声外推~\cite{temme2017error,kandala2019error}；
变分真机实验从分子光谱做到自旋链~\cite{colless2018robust,yu2023simulating}。
我们使用的两个诊断量各有专门文献。各向异性自旋-$1/2$ Heisenberg 链的动力学
自旋结构因子已由 Pereira 等人算出~\cite{pereira2006dynamical}，其动量分辨
响应的静态极限正是我们的 $S(\pi)$ AFM 标记的基础。string 序先为价键相引入
~\cite{dennijs1989preroughening}，后在量子自旋格点中作对称性刻画
~\cite{perezgarcia2008string}，并在低维 Mott 绝缘体中经关联粒子-空穴对观测到
~\cite{endres2011observation}；我们的 string 算符是同一扫参下的可直接测量对应物。
在真机侧，有限连通使 qubit 映射成为首要问题~\cite{li2019tackling}；
我们的稳定子加读出筛选以测量优先回答同一问题，直接排出实验将用的子图名次。
读出缓解覆盖关联展开~\cite{bravyi2021mitigating,nation2021scalable}与无模型
期望值修正~\cite{vandenberg2022modelfree}；我们的逐比特 $2\times2$ 带非负约束
展开按构造属于该张量积家族。有限尺寸能隙标度是自旋链数值的常规做法
~\cite{sandvik2010computational}，Haldane 能隙问题是其经典情形
~\cite{haldane1983nonlinear,affleck1987rigorous,affleck1989quantum}；
我们的 $L=8$ 相划分加能隙标度检查沿用该规范，完整曲线见附录。
本文把相匹配变分制备推广到二聚化 XXZ 链，其中平庸、拓扑与反铁磁三相竞争，
并以可直接测量的标记（$S(\pi)$、string）连同反射不变量
$\tilde{Z}_{\mathcal R}$ 刻画制备态。
我们的方法把相匹配制备推广到三序竞争的单模型，同时从一开始就计入真机约束：
优质比特筛选、读出缓解、标记只限可直接测量的关联子。
制备态能否在真机上复现模拟机趋势，由第~\ref{sec:comparison}~节的比较裁决。

本文组织如下。先引入模型及其对称性（第~\ref{sec:model}~节）并画出有限尺寸相图
（第~\ref{sec:phasediagram}~节，能隙标度见附录），再定义可直接测量的相标记
（第~\ref{sec:markers}~节）、制备变分态（第~\ref{sec:ansatz}~节），把协议搬上真机
（第~\ref{sec:hardware}~节），最后比较真机与模拟机结果（第~\ref{sec:comparison}~节）。

\FloatBarrier
\section{模型的引入与对称性}
\label{sec:model}

我们考虑 $L$ 格点开边界二聚化 XXZ 链，格点编号 $1,\dots,L$。
交替单双键的二聚化结构可追溯到 Su、Schrieffer 和
Heeger~\cite{su1979solitons}。
记第 $m$ 格点上的 Pauli 算符为 $X_m,Y_m,Z_m$，哈密顿量为
\begin{equation}
H(s,\delta)=H_o+H_e,
\end{equation}
其中 $H_o$（$H_e$）耦合奇（偶）键，$s$ 为二聚化参数，$\delta$ 控制各向异性：
\begin{equation}
H_o=(1-s)\sum_{j=1}^{L/2}h_{2j-1,2j},\qquad
H_e=s\sum_{j=1}^{L/2-1}h_{2j,2j+1},
\end{equation}
\begin{equation}
h_{kl}=e^{-\delta}(X_kX_l+Y_kY_l)+e^{+\delta}Z_kZ_l,
\end{equation}
其中 $h_{kl}$ 为双格点键算符，$J_{xx}=e^{-\delta}$，$J_{z}=e^{+\delta}$；
$\delta=0$ 为 Heisenberg 各向同性点，$\delta>0$ 偏向 Ising。
如无特别说明取 $L=8$。

$H$ 与总磁化 $Z_{\rm tot}=\sum_m Z_m$、整体自旋翻转 $\prod_m X_m$ 及全链反射
$R$（$m\leftrightarrow L+1-m$）对易，因此在 $H$ 下演化永不离开初态对称性扇区。
$\delta\ne0$ 时对称性为 $U(1)$（$Z_{\rm tot}$ 守恒）；$\delta=0$ 时扩大为 $SU(2)$。

这里的拓扑相属于一维对称保护家族，其研究始于 Haldane 关于整数自旋反铁磁链
存在有限能隙的猜想~\cite{haldane1983nonlinear,dennijs1989preroughening}，
以及自旋 $1$ 链相图中特殊点上严格价键固体基态的发现~\cite{affleck1987rigorous,dennijs1989preroughening}。
如今这套唯象学已有分类：一维有隙对称相由二阶上同调群即对称性的射影表示
标记~\cite{chen2011classification}，等价地可用矩阵乘积态表述~\cite{schuch2011classifying}，
操作性指纹是纠缠谱简并~\cite{pollmann2010entanglement,li2008entanglement}。
我们的二聚化 XXZ 链在单一可调模型内容纳平庸、拓扑、反铁磁三序竞争，
正因如此才需要一个能同时分辨三者的标记。

为区分各相，我们使用归一化拓扑反映标记，采用 Elben 等人的
部分反射多体不变量~\cite{elben2020manybody}。记基态在中央块 $I$ 上
（$L=8$ 时取格点 $3$--$6$）的约化密度矩阵为 $\rho_I$，
${\mathcal R}_I$ 为关于中央键的镜像交换。以
$Z_{\mathcal R}={\rm Tr}(\rho_I{\mathcal R}_I)$ 及 $I$ 左右两半
对应的量 $\rho_{I_{1,2}}$，有
\begin{equation}
\tilde{Z}_{\mathcal R}=
\frac{Z_{\mathcal R}}
{\sqrt{[{\rm Tr}(\rho_{I_1}^2)+{\rm Tr}(\rho_{I_2}^2)]/2}},
\end{equation}
其中分母为两半的平均纯度。换言之，$\tilde{Z}_{\mathcal R}$ 是以局域混合度
归一化的反射期望，在平庸（拓扑）相趋于 $+1$（$-1$）、磁有序相趋于 $0$，
与整体纠缠水平无关。以此标记画出的有限尺寸相图见第~\ref{sec:phasediagram}~节。

\FloatBarrier
\section{有限尺寸相图}
\label{sec:phasediagram}

$99\times99$ 网格 $s=i/100$、$\delta=j/50$（$i,j=1,\dots,99$）上的归一化标记
见图~\ref{fig:D01}。每格点由存储的基态独立计算，不插值、不平滑；
简并点保留求解器最低本征矢，数据集中携带简并标记。
该图用于定义参数区域、代表点选择和后续变分起点/终点协议。
这里将归一化标记作为相标签约定；边界的有限尺寸漂移在附录中单独分析。

\begin{figure}[htbp]
\centering
\includegraphics[width=\columnwidth]{exp01_D01.pdf}
\caption{$L=8$ 时 $99\times99$ 网格上的归一化标记
$\tilde{Z}_{\mathcal R}(s,\delta)$，逐点计算不插值。
\label{fig:D01}}
\end{figure}

\FloatBarrier
\section{相标记}
\label{sec:markers}

\subsection{结构因子与 AFM 标记}
\label{sec:sq}

为识别反铁磁相，我们使用结构因子
\begin{equation}
S(q)=\frac{1}{L}\sum_{i,j}e^{iq(i-j)}\langle Z_iZ_j\rangle,
\end{equation}
其中 $\langle Z_iZ_j\rangle$ 为基态关联子，本节恒取 $L=8$。
即 $S(q)$ 是 N\'eel 型关联的傅里叶权重，长程反铁磁序表现为 $q=\pi$ 处的
Bragg 型尖峰。三相各一个代表点处的 $S(q)$ 曲线（相标签与坐标见数据集），
共用 $q$ 网格（$99$ 点，含 $q=\pi$），见图~\ref{fig:D03}。
在面向真机的测量协议中，使用 $q=\pi$ 分量作为 AFM 诊断量；
其对不同相的实际区分结果放在结果部分报告，而不在方法部分预先断言。

\begin{figure}[htbp]
\centering
\includegraphics[width=\columnwidth]{exp02_D03.pdf}
\caption{共用 $q$ 网格上三相各一个代表点处的结构因子 $S(q)$；
只有 AFM 点在 $q=\pi$ 处见峰。
\label{fig:D03}}
\end{figure}

\FloatBarrier
\subsection{拓扑热力图}
\label{sec:heatmaps}

同一 $(s,\delta)$ 网格上的四幅热力图见图~\ref{fig:D04}：(a) AFM 结构因子
$S(\pi)$；(b) string 算符，自 den Nijs 和 Rommelse~\cite{dennijs1989preroughening}
以来用于分辨对称保护拓扑序；(c) ZR-like 组合
$Q=\frac{4}{3}+2O_{\rm str}-\frac{1}{6}S(\pi)$，其中 $O_{\rm str}$ 为 string
期望；(d) 归一化标记 $\tilde{Z}_{\mathcal R}$（Ref.~\cite{elben2020manybody}，
此处独立于图~\ref{fig:D01}重算）。四个量在同一参数网格上计算，
以便在结果部分统一比较其诊断作用。
图 (b) 的 string 信号是如下关联子的长程极限：
\begin{equation}
O^z_{\rm str}(i,j)=\Big\langle S^z_i\exp\!\Big(i\pi\!\sum_{k=i+1}^{j-1}S^z_k\Big)S^z_j\Big\rangle,
\end{equation}
其中 $S^z_k$ 为格点 $k$ 的自旋 $z$ 算符：相位 string 探测任何局域两点函数都
看不见的隐藏反铁磁序，即为价键相引入的序参量~\cite{dennijs1989preroughening}，
如今从对称保护的纠缠谱简并理解~\cite{pollmann2010entanglement}。
这类纠缠指纹始于纠缠谱作为拓扑诊断~\cite{li2008entanglement}与纠缠熵的普适常数项
$S=\alpha L-\gamma$~\cite{kitaev2006topological,levin2006detecting}；
而我们的图 (a)(b) 以单比特读出可测的关联子换取了全纠缠分辨。
随机测量如今以很少的设置重构这类量：classical shadow 以
$N\ge{\cal O}(\log M\,\max\|O_i\|^2_{\rm shadow}/\epsilon^2)$ 次测量同时预言
$M$ 个线性性质~\cite{huang2020predicting}，可去随机化为确定性 Pauli
测量~\cite{huang2021derandomization}、向目标哈密顿量偏置~\cite{hadfield2022measurements}，
推广到多 shot 估计~\cite{zhou2023performance}、互无偏基~\cite{wang2023classical}乃至
信道层析~\cite{kunjummen2023shadow}，见综述~\cite{elben2022toolbox}；
二阶 R\'enyi 熵与量子 Fisher 信息已如此测得~\cite{brydges2019probing,vitale2024robust}。
与 $\tilde{Z}_{\mathcal R}$ 不同，图 (a)(b) 只含标准 Pauli 关联子，可直接测量。
由于 $\tilde{Z}_{\mathcal R}$ 不易在真机上直接测量，面向真机的协议使用态制备，
并结合图 (a)(b) 中可直接测量的量。

\begin{figure*}[p]
\centering
\begin{minipage}[b]{0.48\textwidth}
\centering
\includegraphics[width=\textwidth]{exp02_D04a.pdf}\\(a)
\end{minipage}
\hfill
\begin{minipage}[b]{0.48\textwidth}
\centering
\includegraphics[width=\textwidth]{exp02_D04b.pdf}\\(b)
\end{minipage}
\\
\begin{minipage}[b]{0.48\textwidth}
\centering
\includegraphics[width=\textwidth]{exp02_D04c.pdf}\\(c)
\end{minipage}
\hfill
\begin{minipage}[b]{0.48\textwidth}
\centering
\includegraphics[width=\textwidth]{exp02_D04d.pdf}\\(d)
\end{minipage}
\caption{图~\ref{fig:D01} $(s,\delta)$ 网格上的拓扑热力图：(a) $S(\pi)$，
(b) string 算符，(c) $Q$，(d) $\tilde{Z}_{\mathcal R}$。\label{fig:D04}}
\end{figure*}

\FloatBarrier
\section{变分态制备}
\label{sec:ansatz}

\subsection{示意起点与终点}
\label{sec:points}

三对示意起点/终点，每相一对，标在 $\tilde{Z}_{\mathcal R}$ 相图上，
见图~\ref{fig:D05}：平庸相 $(0.01,0.02)\to(0.21,0.42)$，
$\tilde{Z}_{\mathcal R}\approx0.996$；拓扑相 $(0.99,0.02)\to(0.90,0.20)$，
$\tilde{Z}_{\mathcal R}\approx-0.98$；AFM 相 $(0.50,1.98)\to(0.61,1.44)$，
$\tilde{Z}_{\mathcal R}\approx0.009$。
终点是相深处的经典参考目标，待相图落定后选取——变分运行的示范终点，
并非真实演化终点本身。终点置于相深处，使参考路径避开有限尺寸漂移最大的
边界区。起点沿用三色圆点样式；终点标记形状颜色均区别于起点，
箭头给出每对的参考方向。

\begin{figure}[htbp]
\centering
\includegraphics[width=\columnwidth]{exp01_D05.pdf}
\caption{图~\ref{fig:D01} $\tilde{Z}_{\mathcal R}$ 相图上的示意起点/终点对，
每相一对、参考终点在相深处；颜色区分各对，箭头标参考方向。
\label{fig:D05}}
\end{figure}

\subsection{变分电路结构}

我们的变分电路遵循变分本征求解器范式~\cite{peruzzo2014variational}：
制备初态，再作用参数化拟设。经典循环最小化
\begin{equation}
E(\theta)=\langle\psi(\theta)|H|\psi(\theta)\rangle,
\end{equation}
其中 $|\psi(\theta)\rangle$ 为参数 $\theta$ 下的电路输出。
这类变分算法见综述~\cite{cerezo2021variational,tilly2022variational}，
其中收集了代价函数设计、拟设族与最佳实践。
实际瓶颈是优化景观而非表达能力：随机电路形成 $2$-design 时梯度随比特数
指数消失~\cite{mcclean2018barren}；硬件高效电路随深度出现陡峭的可训练性
转变~\cite{campos2021abrupt}；最优参数可以问题规模的逆多项式集中~\cite{akshay2021parameter}；
光滑的哈密顿量变分解跨尺寸迁移且梯度不消失~\cite{mele2022avoiding}；
贫瘠高原区还可藏指数多个陷阱~\cite{nemkov2025barren}。
Krylov 型迭代求解器是到基态的另一条路~\cite{bharti2021iterative}。
硬件高效变分优化最早在超导真机上对小分子与量子磁体实现~\cite{kandala2017hardware}，
自旋链变分模拟随后在云超导设备上推到 $102$ 比特~\cite{yu2023simulating}。
我们的 orbit 拟设与之都不同：以反射对称性绑定参数、初态定纠缠结构，详述如下。
三初态（奇键单态乘积及其对应态，均在 $P=-2$ 扇区，$P$ 标记组合对称性本征值）
配 orbit 拟设，镜面对映键对共享转角（奇键 $O_j\leftrightarrow O_{M+1-j}$，
偶键 $E_j\leftrightarrow E_{M-j}$）。即拟设分两步构造：先以相匹配初态定下
纠缠结构，再给所有键穿上由反射对称性绑定的转角，参数量随键 orbit 数而非键数增长。
每相一个示意图见图~\ref{fig:D06}(a)--(c)：初态 + 所需拟设电路，分别为平庸、
拓扑、AFM 相。
该构造保证试探态永不离开目标对称性扇区；交叠是否足够，由第~\ref{sec:comparison}~节的
真机-模拟机基准裁决。

\begin{figure}[htbp]
\centering
\includegraphics[width=\columnwidth]{exp04_D06a.pdf}
\caption{(a) 平庸初态 + 所需拟设电路。虚线框 $U^{(k)}$ 为第 $k$ 变分层；
$\theta^{(k)}_j$（$\phi^{(k)}_j$）为第 $k$ 层第 $j$ 镜面对映键 orbit 共用的
$XX{+}YY$（$ZZ$）转角（奇键 $O_j\leftrightarrow O_{M+1-j}$，偶键
$E_j\leftrightarrow E_{M-j}$；$\delta=0$ 时 $\theta=\phi$）。偶键子层先作用。
\label{fig:D06}}
\end{figure}

\begin{figure}[htbp]
\centering
\addtocounter{figure}{-1}
\includegraphics[width=\columnwidth]{exp04_D06b.pdf}
\caption{(b) 拓扑初态 + 所需拟设电路，记号同 (a)；奇键子层先作用。}
\end{figure}

\begin{figure}[htbp]
\centering
\addtocounter{figure}{-1}
\includegraphics[width=\columnwidth]{exp04_D06c.pdf}
\caption{(c) AFM 初态 + 所需拟设电路，记号同 (a)；偶键子层先作用。}
\end{figure}

\FloatBarrier
\section{真机实现}
\label{sec:hardware}

\subsection{真机优质比特选择}
\label{sec:premium}

为上真机，我们用两类量筛选候选链（环）：稳定子均值探双比特门质量，
读出保真度探测量质量。理想输出是 cluster 态
$|C\rangle=\prod_{\langle ij\rangle}CZ_{ij}|+\rangle^{\otimes n}$，
其中 $\langle ij\rangle$ 取遍候选子图的边；其稳定子 $S_i$
（开端点二体 $X_0Z_1$、$Z_{n-2}X_{n-1}$，其余三体 $Z_{i-1}X_iZ_{i+1}$）
无误差时期望为 $+1$，故 $\langle S_i\rangle$ 偏离 $1$ 直接反映 $CZ$ 门误差。
每条候选跑 $4$ 个基准电路：\texttt{stab\_g0} 与 \texttt{stab\_g1} 分测偶、奇
指标稳定子（两组反对易，不能同时读出），\texttt{allzero} 与 \texttt{allone}
不含 $CZ$ 门，以 $P(00\cdots0)$ 与 $P(11\cdots11)$ 的平均给出读出保真度
$F_{\rm ro}$。按下式打分排序
\begin{equation}
{\rm score}=0.8\,\bar{S}+0.2\,F_{\rm ro},
\end{equation}
其中 $\bar{S}=n^{-1}\sum_i\langle S_i\rangle$ 为稳定子均值，
双比特门质量权重高于读出。
cluster 态~\cite{raussendorf2001oneway}是其稳定子所测的理想输出；
基于这类态的测量型计算见综述~\cite{briegel2009measurement}。
$8$ 链与 $10$ 环基准电路不同（环多一个闭合 $CZ$、稳定子全为三体），
只各自内部比较，绝不跨几何比较。与通用随机基准不同，该筛选探的正是
真机将要制备的态的稳定子。
真机刻画还有若干互补手段：Pauli-frame 随机化已在超导比特上实现并以门集层析
验证噪声 Pauli 化~\cite{ware2021experimental}；时间关联 dephasing 下随机基准的
分析指出了标准 Markov 衰减模型失效之处~\cite{qi2021randomized}。
基于阴影的哈密顿量测量是省测量的替代方案~\cite{hadfield2022measurements,huang2021derandomization}，
阴影层析还可推广到量子信道~\cite{kunjummen2023shadow}。

具体排名和选中的比特将在 Results and Discussion 中报告。
本方法节只定义筛选协议以及同一几何内部的比较规则。

\subsection{优质比特排名}
\label{sec:rankings}

筛选在百花芯片上报告三组排名：(a) $8$ 链前十；(b) $10$ 环前十；(c) 冠军环内十条 $8$-子链。具体比特表跟随最新 S04；在 D07 图落定前，下方仅保留空位，不做任何排名断言。

\begin{figure}[htbp]
\centering
\fbox{\parbox{0.9\columnwidth}{\centering TODO(GAP-001)：D07a 排名图待出。}}\\(a)
\fbox{\parbox{0.9\columnwidth}{\centering TODO(GAP-001)：D07b 排名图待出。}}\\(b)
\fbox{\parbox{0.9\columnwidth}{\centering TODO(GAP-001)：D07c 排名图待出。}}\\(c)
\caption{优质比特排名占位。\label{fig:D07}}
\end{figure}

\FloatBarrier
\subsection{读出误差缓解}
\label{sec:readout}

读出误差把理想分布变成测量分布：
\begin{equation}
\vec{P}_{\rm meas}=M\vec{P}_{\rm ideal},
\end{equation}
其中 $M$ 为分配矩阵，$M_{ab}=P(a|b)$。
其列由制备标定：制备 $|0\rangle$ 得第一列，制备 $|1\rangle$ 得第二列，
\begin{equation}
M=\begin{bmatrix} P(0|0) & P(0|1) \\ P(1|0) & P(1|1) \end{bmatrix},
\end{equation}
其中对角元为读对概率，非对角元为读错概率。矫正时求逆，
\begin{equation}
\vec{P}_{\rm mit}=M^{-1}\vec{P}_{\rm meas},
\end{equation}
其中 $\vec{P}_{\rm mit}$ 再经负值截零、重归一投影到单纯形。
换言之，原始计数先以标定的混淆矩阵展开，再算观测量，采用
Bravyi 等人的张量积读出缓解~\cite{bravyi2021mitigating}。
假设各比特读出误差不关联，逐比特独立矫正，标定代价每比特 $2$ 个电路，
随比特数线性；真机上按批次重标定防漂移，并关闭平台自带读出矫正。
不关联假设保住线性开销，但吃不下串扰；若存在关联读出误差，缓解后仍有残留。
该张量积求逆处于更大的缓解版图之中。概率误差抵消以准概率基展开含噪门并
重采样~\cite{temme2017error}，现已推广到含中途测量的动态电路~\cite{gupta2024probabilistic}，
并与噪声放缩统一~\cite{mari2021extending}；零噪声外推则在放大的噪声强度下拟合
$E(\lambda)$ 并延拓到 $E(0)$——最早在超导真机上实现~\cite{kandala2019error}——
有优化的 Richardson 协议~\cite{krebsbach2022optimization}、逆电路噪声
估计~\cite{koenig2024inverted}，以及变分循环内以人工放大噪声主动压误差~\cite{li2017efficient}；
但时间关联噪声系统性地更难外推掉~\cite{schultz2022analyzing}。
无矩阵迭代展开把同一分配矩阵思想推到数十比特而不显式构造 $M$~\cite{nation2021scalable}；
我们的逐比特求逆是其最简、抗漂移的端点。

\FloatBarrier
\section{真机-模拟机对比}
\label{sec:comparison}

真机与模拟机置于同一图中比较：模拟机 $p=1,2,3$ 曲线叠放真机最优层 $p^*$
散点，每个 estimator 带标准差误差棒。
即每点给出各自 estimator 的采样不确定度，层间散布可对照 shot 噪声判断，
而不读作物理。真机与模拟机各自内部作 $0$--$1$ 归一化，只比变化趋势，
不比绝对数值。分别归一是刻意的：门误差压制与读出残差对两平台绝对标尺的
平移不同，而 AFM 结构因子与 string 算符随 $s$ 的变化作为共同信号留存。
真机变分实验构成背景：双 transmon 上 VQE 加子空间展开算全 H2
能谱~\cite{colless2018robust}、囚禁离子上量子化学~\cite{hempel2018quantum}、
硬件高效优化六比特分子与磁问题~\cite{kandala2017hardware}、超导十二比特上
Hartree--Fock~\cite{google2020hartree}、云设备上 $102$ 比特自旋链基态~\cite{yu2023simulating}。
我们的协议不同在：以可直接测量的拓扑标记、逐层对照真机与模拟机，而不只比能量。
四图覆盖 $\delta=0$ 与 $\delta=0.85$ 乘两个可直接测量标记；图位预留如下，
趋势论述待 D08 图落定前保持冻结。

\begin{figure*}[htbp]
\centering
\begin{minipage}[b]{0.48\textwidth}
\centering
\fbox{\parbox{0.9\textwidth}{\centering TODO(GAP-002)：D08a 图待出（$\delta=0$ AFM）。}}\\(a)
\end{minipage}
\hfill
\begin{minipage}[b]{0.48\textwidth}
\centering
\fbox{\parbox{0.9\textwidth}{\centering TODO(GAP-002)：D08b 图待出（$\delta=0$ string）。}}\\(b)
\end{minipage}
\\
\begin{minipage}[b]{0.48\textwidth}
\centering
\fbox{\parbox{0.9\textwidth}{\centering TODO(GAP-002)：D08c 图待出（$\delta=0.85$ AFM）。}}\\(c)
\end{minipage}
\hfill
\begin{minipage}[b]{0.48\textwidth}
\centering
\fbox{\parbox{0.9\textwidth}{\centering TODO(GAP-002)：D08d 图待出（$\delta=0.85$ string）。}}\\(d)
\end{minipage}
\caption{真机-模拟机对比四图（占位）：(a) $\delta=0$ AFM 结构因子，
(b) $\delta=0$ string 算符，(c) $\delta=0.85$ AFM 结构因子，
(d) $\delta=0.85$ string 算符。每图将叠放模拟机 $p=1,2,3$ 曲线与真机 $p^*$
散点；归一化见正文。\label{fig:D08}}
\end{figure*}

\FloatBarrier
\section{结论}
\label{sec:conclusion}

我们引入二聚化 XXZ 链，用归一化反射标记 $\tilde{Z}_{\mathcal R}$ 画出
$L=8$ 有限尺寸相图，分辨平庸、拓扑与 AFM 三相。能隙标度确认有限尺寸效应
只动相边界。与需全层析的标记不同，我们的真机协议完全建立在可直接测量的
关联子上：保对称 orbit 拟设从匹配初态出发制备三相，优质比特筛选选出最优
真机子图，读出缓解展开原始计数。
该协议刻意模块化——更好的拟设、更丰富的筛选指标、关联读出模型都可各自升级
而不动比较框架——为在更大格点上追踪拓扑标记提供模块化起点。
本研究的局限有三：相划分全系于 $L=8$、尺寸外推只有能隙标度一路；
读出缓解假设各比特误差不关联；真机-模拟机趋势比较待 D08 图，
在此之前真机的比较一半是预留的协议而非结果。

\FloatBarrier
\appendix

\section{有限尺寸能隙说明}
\label{sec:gap}

取 $P=-2$ 对称性扇区的能隙
\begin{equation}
\Delta_{\rm sec}(L,s)=E_1^{(P=-2)}(L,s)-E_0^{(P=-2)}(L,s),
\end{equation}
其中 $E_0^{(P=-2)}$（$E_1^{(P=-2)}$）为该扇区基态（第一激发态）能量。
$\delta=0$ 时 $L=8,12,16$ 的扇区能隙见图~\ref{fig:D02a}。
临界点附近有限尺寸标度给出 $\Delta_{\rm sec}L\approx A$；
画 $A$ 随 $s$ 变化（图~\ref{fig:D02b}），$L$ 增大时交点稳定在 $s=0.5$。
具体说，相邻尺寸间交点的漂移随 $L$ 增大而缩小，与热力学极限下收敛到
$s=0.5$ 一致。$s=0.5$ 处能隙随 $1/L$ 的变化（$L=8,12,16,20,24$）
见图~\ref{fig:D02c}；过 $L=20,24$ 的线性拟合（系数标于图，
拟合截距小，$b=0.05$）与能隙缩小趋势一致。因此有限尺寸效应只对相边界
有一点影响，不影响相大体结构。
全文基态均由精确对角化（Lanczos）得到，交点分析遵循量子自旋系统
有限尺寸常规做法，见 Ref.~\cite{sandvik2010computational} 讲义。

\begin{figure}[htbp]
\centering
\includegraphics[width=\columnwidth]{exp03_D02a.pdf}
\caption{$\delta=0$ 时 $L=8,12,16$ 的对称性能隙 $\Delta_{\rm sec}$。
\label{fig:D02a}}
\end{figure}

\begin{figure}[htbp]
\centering
\includegraphics[width=\columnwidth]{exp03_D02b.pdf}
\caption{$A=\Delta_{\rm sec}L$ 随 $s$ 变化；$L\to\infty$ 时交点稳定在
$s=0.5$。\label{fig:D02b}}
\end{figure}

\begin{figure}[htbp]
\centering
\includegraphics[width=\columnwidth]{exp03_D02c.pdf}
\caption{$s=0.5$ 处能隙随 $1/L$ 变化，过 $L=20,24$ 线性拟合
（拟合截距小，$b=0.05$）。
\label{fig:D02c}}
\end{figure}


\FloatBarrier
\bibliography{references}

\end{document}
